% Tento soubor nahraďte vlastním souborem s obsahem práce.
%=========================================================================
% Autoři: Michal Bidlo, Bohuslav Křena, Jaroslav Dytrych, Petr Veigend a Adam Herout 2019

% Pro kompilaci po částech (viz projekt.tex), nutno odkomentovat a upravit
%\documentclass[../projekt.tex]{subfiles}
%\begin{document}

\chapter{Úvod}
Podpis je jednou z~nejstarších metod používanou pro ověření totožnosti, určení autorství či udělení právního souhlasu.
V~této práci půjde především o~ověření totožnosti, které spadá pod obor zvaný biometrie.
Biometrie se zabývá analýzou biologických a behaviorálních charakteristik používaných mimo jiné k~autentizaci. 
Zkoumáním pravosti podpisu a jiných textů se zabývá písmoznalectví neboli grafognózie.
U~podpisu lze analyzovat statické a dynamické parametry písma. 

%Čím se tedy zabývá tato práce?
Podpis je v~praxi stále nejpoužívanější způsob autentizace. 
Podaří-li se důvěryhodně podpis napodobit, bude možnost obejít bezpečnostní zabezpečení. 
To představuje velké bezpečnostní riziko a nutnost upravení dosavadních bezpečnostních metod pro autentizaci. 
Bavíme se zde spíše o~digitálním podpisu. 
Na rozdíl od podpisu na papíře digitální podpis obsahuje i dynamické charakteristiky.
Bude tedy o~něco složitější jej zfalšovat.


%TOLE JE SPIS ABSTRAKT
%Bude třeba sestavit snímací zařízení pro získaní potřebných parametrů podpisu.
%Ze statických parametrů budeme snímat tvar i tloušťku čar a celkový vzhled podpisu. 
%Z dynamických parametrů budeme snímat tlak, sklon a průběh tahů. 
%Získané data je potřeba vhodným způsobem zpracovat a provést potřebné výpočty. 
%Upravená data je poté potřeba převést do instrukcí pro mechanickou paži, která je použita k replikaci podpisu. 
%Výsledný podpis bude porovnáván s vlastnoručním podpisem. Cílem je co nejvíce se přiblížit vlastnoručnímu podpisu.

%\section{Představení tématu}
%\section{Cíle práce}

%\section{Struktura práce}

\chapter{Teoretické základy}
Pro plné pochopení problematiky je potřeba představit základy z~několika odvětví. 
V~této části budou shrnuty všechny potřebné informace. 

\section{Pojmy}
\label{sec:pojmy}
\begin{itemize}
  \item{\textbf{Biometrie} - automatizované rozpoznávání lidských jedinců na základě jejich charakteristických anatomických a behaviorálních rysů.} % přepsat - citace biometrie (str. 13)
  \item{\textbf{Biologické autentizační metody} - zkoumají biologické charakteristiky člověka, s~jimiž se narodí. Tyto charakteristiky má každý člověk unikátní a jsou neměnitelné. Patří mezi ně například otisky prstů, skeny sítnice či duhovky.}  
  \item{\textbf{Behaviorální autentizační metody} - porovnávají behaviorální charakteristiky související s~chováním člověka. Lze například porovnávat, jakým způsobem jedinec mluví nebo jak pohybuje očima.}
  \item{\textbf{Statické parametry písma} - jde o~parametry, které nezahrnují informace o~procesu psaní podpisu. Patří mezi ně například tvar a vzhled písma, umístění podpisu na stránce nebo také tloušťka čar.}
  \item{\textbf{Dynamické parametry písma} - tyto parametry se vytahují k~samotnému průběhu psaní podpisu. Například rychlost psaní, akcelerace, tlak nebo průběh tahů.}
\end{itemize}

\section{Definice podpisu a jeho charakteristiky}
Definice podpisu je z~pohledu zákona poněkud problematická. 
Podpis jako takový není v~zákoně nijak přímo definován. 
Obecně je ale bráno za podpis vlastnoruční uvedení jména a příjmení, nebo jiného jedinečného a nezaměnitelného označení. % citace? https://www.fulsoft.cz/33/komentar-zakona-89-2012-sb-obcansky-zakonik-561-uniqueidmRRWSbk196FNf8-jVUh4EtuvCojfP1Dmx8t50X_QR_yU_dtwjPItMA/

.Písmo patří z části ke statickým, z části k dynamickým biometrickým vlastnostem.                                                                                  % |
.Záleží, zda se snímá i průběh psaní nebo pouze se uchovává výsledek psaní.                                                                                        % |
.Je pochopitelné, že dynamické vlastnosti obsahují více informace a lze díky nim lépe rozpoznat neprávoplatného uživatele, pokoušejícího se falzifikovat podipis.  % V
.Nicméně k záznamu dynamických vlastností je třeba tablet a speciální aktivní pero, které ukládá patřičné parametry během psaní.                                   % Citace Biometrie (str 233)

.Charakteristiky podpisu jsou nepřeberné množství informací, podle kterých lze jednotlivé podpisy od sebe odlišit.
.Charakteristiky se dělí na statické a dynamické~\ref{sec:pojmy}. 
.Jsou ovlivňovány jak fyzickým, tak i psychologickým stavem člověka v~době podpisu.
.Mezi charakteristiky patří například celkový vzhled podpisu, tloušťka čar, rychlost psaní, tlak, průběh a pořadí tahů a mnoho dalších, které ale pro nás v~rámci této práce nebudou důležité. %TODO vymyslel jsem si lol

\section{Biometrická autentizace}
Biometrická autentizace probíhá na základě biologických a behaviorálních charakteristik člověka\ref{sec:pojmy}. 
Obě tyto skupiny znaků jsou pro každého člověka unikátní a nezaměnitelné. To je důvodem, proč je lze využít k~autentizaci.
Existuje spousta biometrických metod pro autentizaci, kupříkladu otisk prstu, rozpoznání obličeje, sken duhovky či sítnice, nebo dokonce pomocí používání myši. 
Každá tato metoda má své výhody a nevýhody, zejména se hodnotí přesnost, cena, komfort používání, stálost a velikost vzorku.
Tato práce se zaměřuje na autentizaci podpisem.

Existuje spoustu různých cest, kterými lze biometrické systémy napadnout.                                   % |
V~rámci této práce se zaměříme na seznzor, jenž je první komponentu biometrického systému (obrázek níže).   % V
Tomuto senzoru budeme předkládat falešné biometriky, tedy v~našem případě falešného podpisu.                % parafráze Biometrie (str 78)

\begin{figure}[h]
\centering
\includegraphics[width=0.5\textwidth]{obrazky-figures/placeholder.pdf}
\caption{Biometrický systém}
\label{fig:my-pdf}
\end{figure}

\section{Parametry podpisu}
V~rámci této práce budou snímány následující parametry:
\begin{itemize}
  \item{\textbf{Celkový vzhled podpisu}}
  \item{\textbf{Pohyb pera} ze kterého bude vypočítán celkový průběh podpisu, především pak \textbf{pozice pera} a \textbf{rychlost psaní čar}.} %při psaní pomocí akcelerometru zabudovaném v MPU-6050 
  \item{\textbf{Sklon pera}} %pomocí gyroskopu, který je také součástí MPU-6050.}
  \item{\textbf{Tlak}} %pomocí čtveřice tlakových senzorů Interlink Electronics FSR® 400.}
\end{itemize}
Kombinací těchto nasnímaných parametrů by bylo možné dostatečně věrohodně napodobit původní vlastnoruční podpis. 

\begin{figure}[h]
\centering
\begin{minipage}{0.45\textwidth}
    \centering
    \includegraphics[width=\textwidth]{obrazky-figures/placeholder.pdf}
    \caption{Vzhled dynamického podpisu.}
    \label{fig:first-image}
\end{minipage}\hfill
\begin{minipage}{0.45\textwidth}
    \centering
    \includegraphics[width=\textwidth]{obrazky-figures/placeholder.pdf}
    \caption{Graf dynamických parametrů podpisu v~čase.}
    \label{fig:second-image}
\end{minipage}
\end{figure}

\section{Rozpoznávání falzifikátů}
Rozpoznávání falzifikátů je u~digitálního podpisu automatizované.
Probíhá na základě porovnávání uloženého vzorku v~databázi s~podpisem, kterým se daný člověk pokouší autentizovat. %použito správně slovo autentizce?
Referenční podpis je zprůměrován z~několika vzorků, aby došlo k~potlačení náhodných jevů. %  citace? zav_prace 17183 2.3
Tyto vzorky jsou poskytnuté danou osobou při registraci do systému.

Porovnávací algoritmus extrahuje důležité parametry podpisu, které poté porovnává s~referenčním vzorem. 
Výsledkem takového porovnávání je určité procento shody parametrů.
Následně je vypočteno celkové procento shody podpisů a na jeho základě je autentizace úspěšná či nikoli.
Je velmi důležité, aby algoritmus měl určenou správnou procentuální hodnotu shody.
Pokud by byla špatně nastavena, bylo by příliš mnoho neúspěšných autentizací, které by měly být úspěšné, nebo naopak mnoho úspěšných autentizací, které měly být neúspěšné.

\begin{figure}[h]
\centering
\includegraphics[width=0.5\textwidth]{obrazky-figures/placeholder.pdf}
\caption{Průběh autentizace}
\label{fig:my-pdf}
\end{figure}

\subsection{rozpoznání statických falzifikátů} % preparafrazovat podle Biometrie a identita člověka str 440
Verifikace v~off-line systému se skládá ze tří etap – předzpracování, extrakce biometrických
charakteristik a vyhodnocení. Během předzpracování jsou provedeny některé algoritmy, jako
jsou vyhlazování, zjednodušení, či normalizace. Extrakce probíhá různými způsoby, jedním
z~nich je zjištění geometrických a topologických rysů – uzavřené smyčky, speciální body
nebo místa křížení (Obr. 15). Výsledné vyhodnocení je založeno na vyhodnocování vektorů
charakteristik, může být provedeno porovnáním významových bodů, klasifikátorem sousedů
nebo neuronovou sítí. (1, kap. 13.7) %https://theses.cz/id/odwy80/BP_Vchov.pdf (1) https://aleph.lib.vutbr.cz/F/QR556J3JR4JCE18FMRXDUGQ8FXVA2VL3NX3CYI9PYU2PRQCPTX-02400?func=item-global&doc_library=BUT01&doc_number=000088231&year=&volume=&sub_library=BXFIT

\subsection{rozpoznání dynamických falzifikátů} % preparafrazovat podle Biometrie a identita člověka str 444
V~případě on-line verifikace jsou pomocí HW získány kromě statických i dynamické
charakteristiky – rychlost psaní, tlak pera, pořadí psaní jednotlivých částí apod. Již v~průběhu
psaní podpisu jsou získávány informace o~průměrné a maximální rychlosti psaní, měření
vlastností zakřivení tahů, poměr dlouhých a krátkých tahů, různé délky segmentů podpisu aj..
V~každém časovém okamžiku vzniku podpisu jsou popisovány časovou funkcí další
dynamické charakteristiky – souřadnicové pozice hrotu psacího nástroje, rychlost (v),
zrychlení (a) a tlak hrotu na podložku (p) (Obr. 16). Následně jsou získané charakteristiky
porovnávány se šablonou uloženou v~databázi, získá se tak míra ztotožnění. (1, kap. 13.8) %https://theses.cz/id/odwy80/BP_Vchov.pdf (1) https://aleph.lib.vutbr.cz/F/QR556J3JR4JCE18FMRXDUGQ8FXVA2VL3NX3CYI9PYU2PRQCPTX-02400?func=item-global&doc_library=BUT01&doc_number=000088231&year=&volume=&sub_library=BXFIT


\subsection{Normy identifikace pomocí podpisu}
Identifikace pomocí podpisu podléhá normám ISO/IEC 19794-7:2007 Signature/sign time
series data, ISO/IEC 19794-11:2013 Signature/sign processed dynamic data a ANSI/INCITS
395-2005 Signature/Sign Data.
ISO/IEC 19794-7:2007 popisuje dva formáty výměny dat pro podpis a písmo. První je obecný
formát, druhý je kompaktní, který lze využít pro smart karty. (48) % https://theses.cz/id/odwy80/BP_Vchov.pdf (48) https://www.iso.org/standard/38751.html
ISO/IEC 19794-11:2013 pro účely biometrického porovnání specifikuje obecný formát
výměny dat pro zpracované podpisové/znakové údaje o~chování získané z~časové řady
zachycené pomocí zařízení. (49) % https://theses.cz/id/odwy80/BP_Vchov.pdf (49) https://www.iso.org/standard/51824.html
ANSI/INCITS 395-2005 specifikuje formát výměny dat pro digitalizované písmo nebo podpis
pro účely biometrického zápisu, verifikace nebo identifikace. Tento formát je obecný a může
být využit v~široké škále aplikací využívajících elektronický podpis. (50) % https://theses.cz/id/odwy80/BP_Vchov.pdf (50) https://webstore.ansi.org/standards/incits/ansiincits3952005

\chapter{Analýza současných řešení}
Obecný kec

\section{Současná řešení v~oblasti biometrie}
Podrobný kec

\section{Systém, zařízení a kde se využívají, a jak jsou důvěryhodné}

\section{Hodnocení výhod a nevýhod existujících řešení}

\chapter{Návrh snímacího pera}
\section{Výběr senzorů a technologií}
\*subsection{EPS32}
\*subsection{MPU-6050 (akcelerometr a gyroskop)}
\*subsection{Tlakový senzor Interlink Electronics FSR® 400}
\section{Schéma a popis návrhu}
\section{Možnosti replikace podpisu}

%\chapter{Implementace prototypu}
%\section{Vývoj hardwaru}
%\section{Sbírání dat v~reálném čase}
%\section{Ukládání dat pro následnou replikaci podpisu}

%\chapter{Replikace podpisu}
%\section{Využití robotické ruky}
%\section{Transformace nasnímaných dat na kód pro replikaci}
%\section{Testování a výsledky replikace}

%\chapter{Hodnocení výsledků}
%\section{Analýza dosažených výsledků}
%\section{Porovnání s~očekáváními}
%\section{Diskuze o~spolehlivosti a přesnosti replikace}

%\chapter{Závěr}
%\section{Shrnutí hlavních poznatků}
%\section{Zhodnocení významu práce}
%\section{Budoucí perspektivy}

%\section{Možná vylepšení a rozšíření}
%\*subsection{Návrhy na zlepšení prototypu}
%\*subsection{Možnosti dalšího výzkumu a vývoje}


%===============================================================================

% Pro kompilaci po částech (viz projekt.tex) nutno odkomentovat
%\end{document}
